\xiaosi

本章首先阐述了自动驾驶规划领域的三大核心挑战——计算效率、可解释性合规和长尾泛化，指出了当前端到端大模型在边缘部署和安全认证方面的不足。随后，系统梳理了国内外在端到端规划、大模型驱动规划、扩散模型规划、液态神经网络、神经架构搜索和深度强化学习等方向的研究进展。

通过文献调研发现：（1）NCP以极少参数实现了可解释的自动驾驶控制，但其布线拓扑完全依赖手工设计；（2）NCP在强化学习Q-value估计场景下的适用性尚未被探索；（3）NCP的结构化稀疏连接相比全连接LTC的优势缺乏定量验证。这些研究空白构成了本文的出发点。
